% This is part of Mes notes de mathématique
% Copyright (c) 2008-2020
%   Laurent Claessens
% See the file fdl-1.3.txt for copying conditions.

%+++++++++++++++++++++++++++++++++++++++++++++++++++++++++++++++++++++++++++++++++++++++++++++++++++++++++++++++++++++++++++
\section{Matrices}
%+++++++++++++++++++++++++++++++++++++++++++++++++++++++++++++++++++++++++++++++++++++++++++++++++++++++++++++++++++++++++++

%Les matrices et les applications linéaires sont deux choses différentes. Une application linéaire\footnote{Définition \ref{DEFooULVAooXJuRmr}.} est une application d'un espace vectoriel vers un autre, et une matrice est un simple tableau de nombres sur lesquels nous définissons des opérations, de telle sorte à fournir une structure d'espace vectoriel. Le lien entre ces opérations et les opérations correspondantes sur les applications linéaires sera fait plus tard.

%TODO: fixer les notations d'ensembles M( n x m, K) et matrice identité dans l'index des notations.

%--------------------------------------------------------------------------------------------------------------------------- 
\subsection{Définitions}
%---------------------------------------------------------------------------------------------------------------------------

\begin{definition}
    Soit un anneau \( A\) ainsi que des entiers \( m\), \( n\) strictement positifs. Notons \( I = \{ 1,\ldots, n \}\times \{ 1,\ldots, m \} \). L'ensemble \( \eM(m\times n, A)\) est alors le \( A \)-module produit \( A^I \) et est appelé ensemble des \defe{matrices}{matrice} \(n\times m\) sur \( A \).
\end{definition}
Si \( M\) est une matrice, nous notons \( M_{ij}\) au lieu de \( M_{(i,j)}\) l'élément correspondant à \( (i,j)\).


\begin{definition}
Quelques ensembles de matrices particuliers.
  \begin{enumerate}
  \item Si \( n=m\), alors:
  \begin{itemize}
    \item nous disons que la matrice est \defe{carrée}{carrée!matrice},
    \item nous notons \( \eM(n,A)\) pour \( \eM(n\times n,A)\),
    \item \( n \) est appelée \defe{ordre}{ordre!d'une matrice carrée} de la matrice.
  \end{itemize}
  \item Si \( n = 1 \), alors la matrice est appelée \defe{matrice-ligne}{matrice-ligne}.
    \item Si \( m = 1 \), alors la matrice est appelée \defe{matrice-colonne}{matrice-colonne}.
  \end{enumerate}
\end{definition}

\begin{remark}    \label{isoNatMatricesPuissances}
    Quelques remarques à propos de structures supplémentaires.
\begin{enumerate}
    \item
        Nous utiliserons (presque) tout le temps des matrices à coefficients dans un corps. Il est clair que, si \( \eK \) est un corps (commutatif), alors \( \eM(n\times m,\eK) \) a une structure d'espace vectoriel sur \( \eK \).
     \item
         On note les isomorphismes naturels \( \eM(1\times m,A) \simeq A^m\) et \( \eM(n\times 1,A) \simeq A^n\).       
\end{enumerate}
\end{remark}

\begin{lemmaDef}
    Avec la multiplication
    \begin{equation}
        \begin{aligned}
             \eM(n\times p,\eK)\times \eM(p\times m,\eK)&\to \eM(n\times m,\eK) \\
             (A,B)&\mapsto C = AB \tq C_{ij} = (AB)_{ij}=\sum_{k=1}^pA_{ik}B_{kj},
        \end{aligned}
    \end{equation}
    l'espace \( \eM(n,\eK)\) est une \( \eK\)-algèbre\footnote{Définition \ref{DefAEbnJqI}.}.
\end{lemmaDef}

\begin{remark}  \label{REM_produitMatrices}
    Quelques remarques là-dessus.
    \begin{enumerate}
        \item
            La motivation de cette définition pour le produit apparaîtra plus loin, mais le Frido n'étant pas un livre d'introduction, j'imagine que le lecteur a déjà une idée.
        \item
            Il n'empêche que, si \( A \) est une matrice-ligne, alors \( AB \) le sera aussi; et de manière similaire, si \( B \) est une matrice-colonne, alors \( AB \) aussi.
        \item   \label{REM_produitMatricesiii}
            Allons un peu plus loin, en utilisant les isomorphismes naturels de \ref{isoNatMatricesPuissances}: il y a des matrices-lignes particulières \( A_{(r)} \) telles que \( (A_{(r)})_{1j} = \delta_{rj} \) (vaut \( 1 \) si \( r = j \) et \( 0 \) ailleurs); lorsqu'elles multiplient une matrice \( B \), il vient\footnote{On a laissé un indice \( i \) là, mais on aurait pu le remplacer\dots\ par quoi, en fait?}
            \begin{equation}
            (A_{(r)}B)_{ij} = \sum_{k=1}^p(A_{(r)})_{ik}B_{kj} = \sum_{k=1}^p \delta_{rk} B_{kj} = B_{rj};
            \end{equation}
            en d'autres termes, multplier \( A_{(r)} \) par \( B \) revient à prendre la \( r \)\ieme{} ligne de \( B \).
            
            De la même façon, les matrices-colonnes particulières \( B_{(s)} \) telles que \( (B_{(s)})_{i1} = \delta_{is} \) donnent, pour une matrice \( A \): \begin{equation}
            (AB_{(s)})_{ij} = \sum_{k=1}^pA_{ik}(B_{(s)})_{kj} = \sum_{k=1}^p A_{ik} \delta_{ks} = A_{is};
            \end{equation}
            en d'autres termes, multplier \( A \) par \( B_{(s)} \) revient à prendre la \( s \)\ieme{} colonne de \( A \).
        \item
            Par ailleurs, sur les matrices carrées d'ordre \( n \) fixé, le produit de deux matrices est bien défini. Ainsi, \( \eM(n, A)\) se voit conférer une structure d'anneau, dont le neutre pour la multiplication est la matrice carrée \( \mtu_n\) (notée aussi \( \mtu\) lorsqu'il n'y a pas d'ambiguïté sur la taille), donnée par
\begin{equation}
    \mtu_{ij}=\delta_{ij} = \begin{cases}
        1    &   \text{si } i=j\\
        0    &    \text{sinon.}
    \end{cases}
\end{equation}
Il est vite vu que si \( M\) est une matrice carrée d'ordre \( n \), alors \( M\mtu=\mtu M=M\).
    \end{enumerate}
\end{remark}

\begin{definition}
    Pour un élément \( M\in \eM(n\times m, A)\) nous définissons encore
    \begin{description}
        \item[La transposée] \( M^t \tq M^t_{ij}=M_{ji}\) (et donc, \( M^t\in \eM(m\times n, A)\)\footnote{Avez-vous remarqué l'inversion de \( n \) et \( m \)? Sans doute que oui, mais c'est toujours bien de le préciser.}
        \item[La trace] \( \tr(M)=\sum_i M_{ii}\). \quext{La trace pour des matrices pas carrées, pourquoi pas\dots}
    \end{description}
\end{definition}

    Nous verrons plus loin en \ref{SUBSECooGPXVooEYwIiJ} que la définition de transposée d'une application linéaire n'est pas tout à fait évidente; elle sera la définition \ref{DefooZLPAooKTITdd}. Ici nous avons bien défini la transposée d'une matrice, pas d'une application linéaire.


\begin{lemma}[\cite{MonCerveau}]        \label{LEMooUXDRooWZbMVN} \label{LEMooRSJTooQEoOtN}
    Si \( A\), \( B\) et \( C\) sont des matrices nous avons
    \begin{enumerate}
        \item
            \( (AB)^t=B^tA^t\),
        \item       \label{ITEMooXDYQooAlnArd}
            \( \tr(ABC)=\tr(CAB)\).
    \end{enumerate}
\end{lemma}

\begin{proof}
    La première est un simple calcul :
    \begin{equation}
        (AB)^t_{ij}=(AB)_{ji}=\sum_kA_{jk}B_{ki}=\sum_kA^t_{kj}B^t_{ik}=(B^tA^t)_{ij}.
    \end{equation}
    Pour la seconde :
    \begin{equation}
        \tr(ABC)=\sum_{ikl}A_{ik}B_{kl}C_{li}=\sum_{ikl}C_{li}A_{ik}B_{kl}=\sum_l(CAB)_{ll}=\tr(CAB).
    \end{equation}
\end{proof}

\begin{normaltext}
    La seconde égalité est importante et est nommée \defe{invariance cyclique}{invariance cyclique!trace} de la trace. Elle sert entre autres nombreuses choses à prouver que la trace d'une matrice d'une application linéaire ne dépend pas de la base choisie. Ce sera la proposition \ref{PROPooRMYQooWkEpJJ}.
\end{normaltext}

\begin{lemma}       \label{LEMooLXAHooPRyHaF}
    Soient des matrices \( M, N\in \eM(n \times m,A)\). Si pour tout \( x \in A^n\) et tout \( y \in A^m\) nous avons
    \begin{equation}
        \sum_{ij}M_{ij}x_iy_j=\sum_{ij}N_{ij}x_iy_j
    \end{equation}
    alors \( M=N\).
\end{lemma}

\begin{proof}
    Il suffit de calculer, pour des matrices \( X \in \in \eM(1 \times n,A) \) et \( Y \in \eM(m \times 1,A) \) bien choisies, les produits \( XAY \) et \( XBY \), qui sont supposés égaux.\footnote{Cela, c'est grâce à l'isomorphisme naturel de la remarque~\ref{isoNatMatricesPuissances}.}
    
    Ici, "bien choisies" signifie en réalité les matrices particulières de la remarque~\ref{REM_produitMatrices}~\ref{REM_produitMatricesiii}. En effet, soit \( r \in \{1,\dots\, n\} \) et \(s \in \{1,\dots\ , m\} \). Avec les mêmes notations que là-bas, on prend \( X = X_{(r)} \) et \( Y = Y_{(s)} \). Pour toute matrice \( M \), la matrice \( XM \) est la \(r\)\ieme{} ligne de \( M \); puis, \( XMY = (XM) Y \) est la \(s\)\ieme{} colonne de \( XM \), donc la la \(s\)\ieme{} colonne,  \(r\)\ieme{} ligne de \( M \), donc en réalité \( M_{rs} \).
    
    Il s'ensuit que, pour tout \( r \in \{1,\dots\, n\} \) et tout \(s \in \{1,\dots\ , m\} \), \( A_{rs} = B_{rs} \); donc  \( A = B \).
\end{proof}


%---------------------------------------------------------------------------------------------------------------------------
\subsection{Transvections, dilatations, permutations}
%---------------------------------------------------------------------------------------------------------------------------

Dans n'importe quel espace de matrices, nous nommons \( E_{ij}\) la matrice remplie de zéros sauf à la case \( ij\) qui vaut \( 1\). Autrement dit
\begin{equation}
    (E_{ij})_{kl}=\delta_{ik}\delta_{jl}.
\end{equation}

\begin{lemma}   \label{LEM_multiplicationMatriceEij}
  Si \( M \) est une matrice, alors:
  \begin{enumerate}
      \item
          \( E_{ij}M \) est une matrice partout nulle, à l'exception de sa \( i \)\ieme{} ligne, qui contient la \( j \)\ieme{} ligne de \( M \);
      \item
          \( ME_{ij} \) est une matrice partout nulle, à l'exception de sa \( j \)\ieme{} colonne, qui contient la \( i \)\ieme{} colonne de \( M \).
  \end{enumerate}
\end{lemma}

\begin{proof}
    Calculons la composante \( kl\) de la matrice \( E_{ij}M\) :
    \begin{subequations}
        \begin{align}
            (E_{ij}M)_{kl}&=\sum_m(E_{ij})_{km}M_{ml}\\
            &=\sum_m\delta_{ik}\delta_{jm}M_{ml}\\
            &=\delta_{ik}M_{jl}\\
            &=
            \begin{cases}
            M_{jl} &\text{si } i=k\\
            0 & \text{sinon}.
            \end{cases}
        \end{align}
    \end{subequations}
    C'est donc la matrice pleine de zéros, sauf la ligne \( i\) qui est donnée par la ligne \( j\) de \( M\).

    De la même manière,
    \begin{subequations}
        \begin{align}
            (ME_{ij})_{kl}&=\sum_mM_{km}(E_{ij})_{ml}\\
            &=\sum_mM_{km}\delta_{im}\delta_{jl}\\
            &=M_{ki}\delta_{jl}\\
            &=
            \begin{cases}
            M_{ki} &\text{si } j=l\\
            0 & \text{sinon},
            \end{cases}
        \end{align}
    \end{subequations}
    donc seule la \( j\) \ieme{} colonne de \(ME_{ij}\) contient (éventuellement) des éléments non-nuls, et c'est précisément la \( i \)\ieme{} colonne de \( M \)
\end{proof}

\begin{definition}
    Une \defe{matrice de transvection}{transvection (matrice)}\index{matrice!de transvection} est une matrice de la forme
    \begin{equation}
        T_{ij}(\lambda)=\mtu+\lambda E_{ij}
    \end{equation}
    avec \( i\neq j\).

    Une \defe{matrice de dilatation}{matrice!de dilatation}\index{dilatation (matrice)} est une matrice de la forme
    \begin{equation}
        D_i(\lambda)=\mtu+(\lambda-1)E_{ii}.
    \end{equation}
    Ici le \( (\lambda-1)\) sert à avoir \( \lambda\) et non \( 1+\lambda\). C'est donc une matrice qui dilate d'un facteur \( \lambda\) la direction \( i\) tout en laissant le reste inchangé.

    Si \( \sigma\) est une permutation\footnote{C'est un élément du groupe symétrique \( S_n\), voir définition~\ref{DEFooJNPIooMuzIXd}.} alors la \defe{matrice de permutation}{matrice!de permutation}\index{permutation!matrice} associée est la matrice d'entrées
    \begin{equation}
        (P_{\sigma})_{ij}=\delta_{i\sigma(j)};
    \end{equation}
    en particulier, pour \( \sigma = (k,l) \), on notera simplement \( P_{kl} \) pour \( P_{\sigma} \).
\end{definition}

\begin{lemma}   \label{LemyrAXQs}
\begin{enumerate}
  \item
    La matrice \( T_{ij}(\lambda)A=(\mtu+\lambda E_{ij})A\) est la matrice \( A\) à qui on a ajouté, à la \( i \)\ieme{} ligne, \( \lambda \) fois la \( j \) \ieme{} ligne; en d'autres termes, c'est \(A \) à qui on a fait la substitution
    \begin{equation}
        L_i\leftarrow L_i+\lambda L_j.
    \end{equation}
    
  \item
    La matrice \( AT_{ij}(\lambda)\)  est la matrice \( A\) à qui on a ajouté, à la \( j \)\ieme{} colonne, \( \lambda \) fois la \( i \) \ieme{} colonne; en d'autres termes, c'est \(A \) à qui on a fait la substitution
    \begin{equation}
        C_j\leftarrow C_j+\lambda C_i.
    \end{equation}

  \item
    La matrice \( D_{i}(\lambda)A\) est la matrice \( A\) où la \( i\)\ieme{} ligne a été multipliée par \( \lambda \); en d'autres termes, c'est \( A \) à qui on a fait la substitution
    \begin{equation}
        L_i\leftarrow \lambda L_i.
    \end{equation}

  \item
    La matrice \( AD_{i}(\lambda)\) est la matrice \( A\) où la \( i\)\ieme{} colonne a été multipliée par \( \lambda \); en d'autres termes, c'est \( A \) à qui on a fait la substitution
    \begin{equation}
        C_i\leftarrow \lambda C_i.
    \end{equation}
    
  \item
    La matrice \( P_{\sigma}A\) est la matrice \( A\) dans laquelle nous avons permuté les lignes avec \( \sigma^{-1}\). En particulier, quand \( \sigma = (i,j) \), \( P_{\sigma}A =  P_{ij}A\) est la matrice \( A \) où les lignes \( i \) et \( j \) ont été inversées, autrement dit, c'est \( A \) qui a subi
     \begin{equation}
        L_i\leftrightarrow L_j.
    \end{equation} 
           
  \item
    La matrice \( AP_{\sigma}\) est la matrice \( A\) dans laquelle nous avons permuté les colonnes avec \( \sigma\). En particulier, quand \( \sigma = (i,j) \), \(AP_{\sigma} = AP_{ij}\) est la matrice \( A \) où les colonnes \( i \) et \( j \) ont été inversées, autrement dit, c'est \( A \) qui a subi
     \begin{equation}
        C_i\leftrightarrow C_j.
    \end{equation} 
 \end{enumerate}
\end{lemma}

\begin{proof}
\begin{subproof}
  \item[Pour les transvections]
    Par distributivité de la multplication de l'addition, on a
    \begin{equation}
        T_{ij}(\lambda)A = \mtu A+\lambda E_{ij}A = A+\lambda E_{ij}A;
    \end{equation}
    on utilise alors le lemme~\ref{LEM_multiplicationMatriceEij} qui implique que:
    \begin{itemize}
    \item
      les lignes autres que \( i \) ne changent pas;
    \item
      à la ligne \(i \), on ajoute \( \lambda\) fois la ligne \( j\).
    \end{itemize}

    En ce qui concerne l'autre assertion sur les transvections, le calcul est similaire:
    \begin{equation}
        AT_{ij}(\lambda) = A+\lambda A E_{ij};
    \end{equation}
    le lemme~\ref{LEM_multiplicationMatriceEij} implique que les colonnes autres que la \( j\)\ieme{} ne changent pas, tandis que la \( j\)\ieme{} prend en plus \( \lambda\) fois la colonne \( i\).
  
  \item[Pour les dilatations]
    Le travail est similaire à celui des transvections. D'une part, on a
    \begin{equation}
        D_i(\lambda)A=\mtu A+(\lambda-1)E_{ii} A =  A+(\lambda-1)E_{ii} A,
    \end{equation}
    et le lemme~\ref{LEM_multiplicationMatriceEij} implique que les lignes autres que la \( i\)\ieme{} ne changent pas, tandis que la \( i \)\ieme{} prend en plus \( \lambda - 1\) fois la même ligne \( i\); or, \( 1 + (\lambda - 1) = \lambda\), si bien que, en réalité, la \( i \)\ieme{} ligne est multipliée par \( \lambda \).
    
    D'autre part, 
    \begin{equation}
        AD_i(\lambda)= A+(\lambda-1)A E_{ii},
    \end{equation}
    et toujours avec le lemme~\ref{LEM_multiplicationMatriceEij}, les colonnes autres que la \( i\)\ieme{} ne changent pas, tandis que la \( i \)\ieme{} est multipliée par \( \lambda \) (pour la même raison que ci-dessus).
    
  \item[Pour les permutations] Nous avons
    \begin{equation}
        (P_{\sigma}A)_{kl}=\sum_m\delta_{k\sigma(m)}A_{ml}=\sum_m\delta_{\sigma^{-1}(k)m}A_{ml}=A_{\sigma^{-1}(k)l}
    \end{equation}
    (remarquons que \(k = \sigma(m) \) si et seulement si \(\sigma^{-1}(k)  = m\) ). Il s'ensuit que la ligne \( k \) de \(P_{\sigma}A \) est la ligne  \(\sigma^{-1}(k) \) de \(A \).
    
    De manière similaire, mais plus simple,
    \begin{equation}
        (AP_{\sigma})_{kl}=  \sum_m A_{km} \delta_{m\sigma(l)} = A_{k\sigma(l)},
    \end{equation}
    donc la colonne \( l \) de \(AP_{\sigma}\) est la colonne \( \sigma(l) \) de \( A \).
    
    Dans le cas d'une transposition\footnote{Définition~\ref{DEFooXNAFooGTbTTJ}}, disons \( \sigma = (i,j) \), on a:
    \begin{itemize}
    \item
      pour \( P_{\sigma}A \), des lignes inchangées sauf les lignes \( i \) et \(j \) qui sont inversées;
    \item
      pour \( AP_{\sigma} \), des colonnes inchangées sauf les lignes \( i \) et \(j \) qui sont inversées.
    \end{itemize}
\end{subproof}      
\end{proof}


%---------------------------------------------------------------------------------------------------------------------------
\subsection{Déterminant}
%---------------------------------------------------------------------------------------------------------------------------

\begin{definition}      \label{DEFooYCKRooTrajdP}
    Si \( A\in\eM(n,\eK)\) nous définissons le \defe{déterminant}{déterminant!matrice} de \( A\) par la formule
    \begin{equation}
        \det(A)=\sum_{\sigma\in S_n}\epsilon(\sigma)\prod_{i=1}^nA_{i\sigma(i)}
    \end{equation}
    où la somme est effectuée sur tous les éléments du groupe symétrique\footnote{Pour le groupe symétrique, c'est la définition \ref{DEFooJNPIooMuzIXd}, le fait que ce soit un groupe fini est le lemme \ref{LEMooSGWKooKFIDyT}, et pour la somme sur un groupe fini c'est la définition \ref{DEFooLNEXooYMQjRo}.} \( S_n\) et où \(\epsilon(\sigma)\) représente la signature\footnote{Définition~\ref{DEFooWPYSooPWuwWO}.} de la permutation \( \sigma\).
\end{definition}
En se souvenant que \( | S_n |=n!\), nous sommes frappés de stupeur devant le fait que le nombre de termes dans la somme croît de façon factorielle (c'est plus qu'exponentiel, pour info) en la taille de la matrice. Cette formule est donc sans espoir pour une matrice plus grande que \( 3\times 3\) ou à la rigueur \( 4\times 4\) à la main. À l'ordinateur, il est possible de monter plus haut, mais pas tellement.

%///////////////////////////////////////////////////////////////////////////////////////////////////////////////////////////
\subsection{Déterminant en petite dimension}
%///////////////////////////////////////////////////////////////////////////////////////////////////////////////////////////

En dimension deux, le déterminant de la matrice
    $\begin{pmatrix}
        a    &   b    \\
        c    &   d
    \end{pmatrix}$
est le nombre
\begin{equation}        \label{EQooQRGVooChwRMd}
     \det\begin{pmatrix}
         a   &   b    \\
         c   &   d
     \end{pmatrix}=\begin{vmatrix}
          a  &   b    \\
        c    &   d
    \end{vmatrix}=ad-cb.
\end{equation}
Ce nombre détermine entre autres le nombre de solutions que va avoir le système d'équations linéaires associé à la matrice.

Pour une matrice $3\times 3$, nous avons le même concept, mais un peu plus compliqué; nous avons la formule
\begin{equation}
    \det
    \begin{pmatrix}
        a_{11}    &   a_{12}    &   a_{13}    \\
        a_{21}    &   a_{22}    &   a_{23}    \\
        a_{31}    &   a_{32}    &   a_{33}
    \end{pmatrix}
    =
    \begin{vmatrix}
        a_{11}    &   a_{12}    &   a_{13}    \\
        a_{21}    &   a_{22}    &   a_{23}    \\
        a_{31}    &   a_{32}    &   a_{33}
    \end{vmatrix}=
    a_{11}\begin{vmatrix}
        a_{22}  &   a_{23}    \\
        a_{32}    &   a_{33}
    \end{vmatrix}-
    a_{12}\begin{vmatrix}
        a_{21}  &   a_{23}    \\
        a_{31}    &   a_{33}
    \end{vmatrix}+
    a_{13}\begin{vmatrix}
        a_{21}  &   a_{22}    \\
        a_{31}    &   a_{32}
    \end{vmatrix}.
\end{equation}

%---------------------------------------------------------------------------------------------------------------------------
\subsection{Mineur, rang}
%---------------------------------------------------------------------------------------------------------------------------

Pour la définition du rang d'une matrice, nous en donnons une qui est clairement inspirée de l'application linéaire associée.
\begin{definition}[\cite{ooKBOMooSkKHvu}]       \label{DEFooCSGXooFRzLRj}
    Le \defe{rang}{rang d'une matrice} d'une matrice de \( \eM(n,\eK)\) est la dimension de la partie de \( \eK^n\) engendrée par ses colonnes.
\end{definition}

Il est possible d'exprimer le rang d'une matrice de façon plus «intrinsèque» via le concept de mineur.
\begin{definition}[\cite{ooLFZTooJWJUed}]
    Les mineurs d'une matrice sont les déterminants de ses sous-matrices carrées.
\end{definition}
Dans la suite nous désignerons souvent par le mot «mineur» la sous-matrice carrée elle-même au lieu de son déterminant.

\begin{proposition}      \label{DEFooVVBYooJbliTi}
    Le rang d'une matrice est la taille de son plus grand mineur non nul.
\end{proposition}


%--------------------------------------------------------------------------------------------------------------------------- 
\subsection{Manipulations de lignes et de colonnes}
%---------------------------------------------------------------------------------------------------------------------------

Nous voudrions savoir ce qu'il se passe avec le déterminant d'une matrice lorsque nous substituons à une ligne ou une colonne une combinaison des autres lignes et colonnes. Lorsque une matrice est donnée, nous notons \( C_j\) sa \( j\)\ieme colonne.

\begin{lemma}[\cite{MonCerveau,ooKYTYooJlzZMp}]        \label{LEMooCEQYooYAbctZ}
    Si \( A\) est une matrice, alors \( \det(A)=\det(A^t)\).
\end{lemma}

\begin{proof}
    Nous commençons par écrire la définition du déterminant :
    \begin{equation}
        \det(A^t)=\sum_{\sigma\in S_n}\epsilon(\sigma)\prod_{i=1}^n(A^t)_{i,\sigma(i)}=\sum_{\sigma}\epsilon(\epsilon)\prod_iA_{\sigma(i),i}.
    \end{equation}
    Pour chaque \( \sigma\) séparément, nous utilisons la proposition \ref{PROPooQMUDooQQVRIe} pour ré-indexer le produit :
    \begin{equation}
        \prod_i A_{\sigma(i),i}=\prod_iA_{i,\sigma^{-1}(i)}.
    \end{equation}
    Nous profitons du fait que l'application \( \varphi\colon S_n\to S_n\) donnée par \( \varphi(\sigma)=\sigma^{-1}\) soit une permutation de \( S_n\) pour appliquer la définition \ref{DEFooLNEXooYMQjRo} et faire la somme sur \( \sigma^{-1}\) :
    \begin{equation}
        \det(A^t)=\sum_{\sigma}\epsilon(\sigma)\prod_iA_{i,\sigma^{-1}(i)}=\sum_{\sigma}\epsilon(\sigma^{-1})\prod_iA_{i,\sigma(i)}=\det(A)
    \end{equation}
    où nous avons utilisé le fait que \(\epsilon(\sigma^{-1})=\epsilon(\sigma)\) (corolaire \ref{CORooZLUKooBOhUPG}).
\end{proof}

Le fait que \( \det(A)=\det(A^t)\) permet, dans toutes les propositions du type «ce qui arrive au déterminant si on change telle ligne ou colonnes» de ne donner qu'une preuve pour la partie «ligne» et déduire automatiquement le cas «colonne». Le lemme suivant donne un exemple d'utilisation.

\begin{lemma}[\cite{MonCerveau}]        \label{LEMooWMQWooGWFlmC}
    Soit une matrice \( A\). Nous considérons la matrice \( B\) obtenue à partir de \( A\) par la permutation de lignes \( L_k\leftrightarrow L_l\) ainsi que la matrice \( C\) obtenue à partir de \( A^t\) par la permutation de colonnes \( C_k\leftrightarrow C_l\).
    
    Alors \( C^t=B\).
\end{lemma}

\begin{proof}
    Calculons les éléments de matrice de \( C\) :
    \begin{equation}
        C_{ij}=\begin{cases}
            (A^t)_{ij}    &   \text{si }  j\neq k, j\neq l\\
            (A^t)_{ik}    &   \text{si } j=l\\
            (A^t)_{il}    &    \text{si }j=k
        \end{cases}=
        \begin{cases}
            A_{ji}    &   \text{si }  j\neq k, j\neq l\\
            A_{ki}    &   \text{si } j=l\\
            A_{li}    &    \text{si }j=k.
        \end{cases}
    \end{equation}
    Ensuite nous prouvons que \( C^t=B\) en écrivant les éléments de \( C^t\) :
    \begin{equation}
        (C^t)_{ij}=C_{ji}=\begin{cases}
            A_{ij}    &   \text{si } i\neq k, i\neq l\\
            A_{kj}    &   \text{si } i=l\\
            A_{lj}    &    \text{si }i=k.
        \end{cases}
    \end{equation}
    Cette dernière expression est la matrice \( A\) après permutation des lignes \( L_k\leftrightarrow L_l\), c'est-à-dire a matrice \( B\).
\end{proof}

Nous rappelons que \( \mtu \) est la matrice «identité», c'est-à-dire celle dont les entrées sont toutes nulles à l'exception de celles pour lesquelles les indices de lignes et de colonnes sont les mêmes; autrement dit, \( \mtu_{ik} = \delta_{ik}\).  Nous écrivons également \( E_{ij}\) la matrice contenant de zéros partout sauf en \( (i,j)\) où elle a un \( 1\), c'est-à-dire
\begin{equation}
    (E_{ij})_{kl}=\delta_{ik}\delta_{jl}.
\end{equation}

\begin{proposition}[Permuter des lignes ou des colonnes \( L_k\leftrightarrow L_l\)\cite{ooKBOMooSkKHvu,MonCerveau}]    \label{PROPooFQRDooRPfuxk}
    Soient une matrice \( A\in \eM(n,\eK)\), deux entiers naturels \( k\neq l\) inférieurs ou égaux à \( n\). 
    \begin{enumerate}
        \item   \label{ITEMooAIHWooHXzeys}
            Si \( B\) est la matrice obtenue à partir de \( A\) en permutant deux lignes ou deux colonnes, alors
            \begin{equation}
                \det(A)=-\det(B).
            \end{equation}
        \item  \label{ITEMooDNHWooOMgmxa}
            Si \( B\) est la matrice obtenue à partir de \( A\) par la permutation de lignes \( L_k\leftrightarrow L_l\). Alors
            \begin{equation}
                B=P_{kl}A
            \end{equation}
            avec \( P_{kl}=\mtu+E_{kl}+E_{lk}-E_{kk}-E_{ll}\) (autrement dit : la matrice \( P_{kl}\) est une matrice de permutations de lignes).
        \item \label{ITEMooSHRQooQrqVdO}
            La matrice \( P_{kl}\) vérifie \( \det(P_{kl})=-1\)
        \item       \label{ITEMooQXSEooMWiKbL}
            Nous avons 
            \begin{equation}
                \det(P_{kl}A)=\det(P_{kl})\det(A).
            \end{equation}
    \end{enumerate}
\end{proposition}

\begin{proof}
    Point par point
    \begin{subproof}
    \item[\ref{ITEMooAIHWooHXzeys} pour les colonnes]

    Soient \( k\) et \( l\) fixés, et considérons la permutation des colonnes \( C_k\) et \( C_l\). Nous notons \( \alpha\) la permutation \( (k, l)\) dans \( S_n\) (groupe symétrique, définition \ref{DEFooJNPIooMuzIXd}). Nous avons
    \begin{equation}
        B_{ij}=A_{i \alpha(j)},
    \end{equation}
    ou encore : \( A_{ij}=B_{i\alpha(j)}\). Par définition,
    \begin{equation}
        \det(A)=\sum_{\sigma\in S_{n}}\epsilon(\sigma)\prod_{i=1}^nA_{i\sigma(i)}
    \end{equation}
    C'est le moment d'utiliser la proposition \ref{PROPooWJQQooFINSEc} à propos de somme sur des groupes avec \( G=S_n\), \( h=\alpha\) et 
    \begin{equation}
        f(\sigma)=\epsilon(\sigma)\prod_iA_{i,\sigma(i)}.
    \end{equation}
    Nous savons que \( \epsilon(\alpha)=-1\) et que \( \epsilon\) est un homomorphisme par la proposition \ref{ProphIuJrC}\ref{ITEMooBQKUooFTkvSu}, donc
    \begin{equation}
        f(\alpha \sigma)=\epsilon(\alpha\sigma)\prod_iA_{i,(\alpha\sigma)(i)}=-\epsilon(\sigma)\prod_iB_{i,\sigma(i)}.
    \end{equation}
    Avec ça, nous concluons :
    \begin{equation}
        \det(A)=\sum_{\sigma\in S_n}f(\sigma)=\sum_{\sigma}f(\alpha \sigma)=-\sum_{\sigma\in S_n}\epsilon(\sigma)\prod_{i=1}^nB_{i\sigma(i)}=-\det(B).
    \end{equation}
    \item[\ref{ITEMooAIHWooHXzeys} pour les lignes]

    Que se passe-t-il si nous permutons les lignes \( L_k\) et \( L_{l}\) ? Si nous notons \( B'\) la matrice obtenue à partir de \( A\) par la permutation de lignes \( L_k\leftrightarrow L_l\), et \( C\) celle obtenue de \( A^t\) après permutation de colonnes \( C_k\leftrightarrow C_l\) alors le lemme \ref{LEMooWMQWooGWFlmC} nous dit que \( C^t=B'\). En utilisant le lemme \ref{LEMooCEQYooYAbctZ} sur le déterminant de la transposée,
    \begin{equation}
        \det(B')=\det(C^t)=\det(C)=-\det(A^t)=-\det(A).
    \end{equation}
    Voila qui prouve le résultat pour les permutations de lignes.
        
\item[\ref{ITEMooDNHWooOMgmxa}]
    Si \( k=l\), il n'y a pas de permutations, et il est vite vu que la matrice \( P_{kl}\) est l'identité parce qu'il y a quatre fois le terme \( E_{kk}\). Nous supposons donc que \( k\neq l\); en particulier \( \delta_{kl}=0\).

    Il s'agit surtout d'un beau calcul :
    \begin{subequations}
        \begin{align}
            (P_{kl}A)_{ij}=\sum_{m}(P_{kl})_{im}A_{mj}&=A_{ij}+\sum_m(\delta_{ki}\delta_{lm}+\delta_{li}\delta_{lm}-\delta_{ki}\delta_{km}-\delta_{li}\delta_{lm})A_{mj}\\
            &=A_{ij}+\delta_{ki}A_{lj}+\delta_{li}A_{kj}-\delta_{ki}A_{kj}-\delta_{li}A_{lj}\\
            &=A_{ij}+\delta_{ki}(A_{lj} - A_{kj})
            +\delta_{li}(A_{kj}-A_{lj}).
        \end{align}
    \end{subequations}
    Si \( i\neq j\) et \( i\neq l\), alors \( (P_{kl}A)_{ij}=A_{ij}\). Si \( i=k\), alors 
    \begin{equation}
        (P_{kl}A)_{kj}=A_{kj}+A_{lj}-A_{kj}=A_{lj},
    \end{equation}
    c'est-à-dire que la \( k\)\ieme ligne de \( SA\) est la \( l\)\ieme ligne de \( A\).

    Avec \( i=l\) nous obtenons la \( k\)\ieme ligne de \( A\).

    Tout cela montre que \( P_{kl}A\) est la matrice \( A\) dans laquelle les lignes \( k\) et \( l\) ont été inversées, c'est-à-dire \( P_{kl}A=B\).

\item[\ref{ITEMooSHRQooQrqVdO}]
            Pour éviter la multiplication des indices, on note \( S = P_{kl}\). En utilisant la définition du déterminant,
            \begin{subequations}
                \begin{align}
                    \det(S)&=\sum_{\sigma\in S_n}\epsilon(\sigma)\prod_{i=1}^nS_{i\sigma(i)}\\
                    &=\sum_{\sigma}\epsilon(\sigma)\prod_i\big( \delta_{i\sigma(i)}+\delta_{ki}\delta_{l\sigma(i)}+\delta_{li}\delta_{k\sigma(i)}-\delta_{ki}\delta_{k\sigma(i)}-\delta_{li}\delta_{l\sigma(i)} \big).
                \end{align}
            \end{subequations}
            Nous utilisons l'associativité et la commutativité du produit pour séparer les facteurs \( i=k\) et \( i=l\) des autres :
            \begin{equation}
                \det(S)=\sum_{\sigma}\epsilon(\sigma)\prod_{\substack{i\neq k\\i\neq l}}\delta_{i\sigma(i)}(\delta_{k\sigma(k)}+\delta_{l\sigma(k)}-\delta_{k\sigma(k)})(\delta_{l\sigma(l)}+\delta_{k\sigma(l)}-\delta_{l\sigma(l)}).
            \end{equation}
            À cause des facteurs \( i\neq k\) et \( i\neq l\), les \( \sigma\) pour lesquels le tout n'est pas nul doivent vérifier \( \delta_{i\sigma(i)}=1\) pour tout \( i\) différent de \( k\) et \( l\). Les deux seuls sont donc \( \sigma=\id\) et la permutation \( \sigma=(k,l)\). Pour \( \sigma=\id\), nous avons
            \begin{equation}
                \prod_{\substack{i\neq k\\i\neq l}}\delta_{ii}(\delta_{kk}+\delta_{lk}-\delta_{kk})(\delta_{ll}+\delta_{kl}-\delta_{ll})=0.
            \end{equation}
            Dernier espoir : \( \sigma=(k,l)\). Pour ce terme nous avons \( \epsilon(\sigma)=-1\) et
            \begin{equation}
                \prod_{\substack{i\neq k\\i\neq l}}\delta_{ii}(\delta_{kl}+\delta_{ll}-\delta_{kl})(\delta_{lk}+\delta_{kk}-\delta_{lk})=1.
            \end{equation}
            Au final dans \( \det(S)\) il n'y a de non nul que le terme \( \sigma=(k,l)\) et il vaut \( -1\). Donc
            \begin{equation}
                \det(S) = \det(P_{kl}) = -1.
            \end{equation}
        \item[\ref{ITEMooQXSEooMWiKbL}]
            Il s'agit de mettre bout à bout les points déjà prouvés :
            \begin{equation}
                \det(P_{kl}A)=-\det(A)=\det(P_{kl})\det(A).
            \end{equation}
    \end{subproof}
\end{proof}

\begin{corollary}[\cite{ooKBOMooSkKHvu}]        \label{CORooAZFCooSYINvBl}
    Soit une matrice \( A\in \eM(n,\eK)\). Si deux lignes ou deux colonnes de \( A\) sont égales, alors \( \det(A)=0\).
\end{corollary}

\begin{proof}
    Si deux colonnes sont égales, la matrice ne change pas lorsqu'on les permute, alors que le déterminant change de signes. La seule possibilité est que \( \det(A)=-\det(A)\), ce qui signifie que \( \det(A)=0\).
\end{proof}

Notons que si pour \( k\neq l\) nous avons \( C_k=\lambda C_l\), alors nous avons aussi \( \det(A)=0\). Nous le prouverons plus tard.

La réciproque n'est pas vraie : il existe des matrices dont le déterminant est nul et dont aucune entrée n'est nulle. Par exempe
\begin{equation}
    \begin{pmatrix}
        1    &   2    \\ 
        1    &   2    
    \end{pmatrix}.
\end{equation}

\begin{probleme}
Cet exemple n'est pas à sa place là. Peut-être 1,2 // 4,8 est mieux? Ça dépend de ce qu'on veut illustrer\dots
\end{probleme}


\begin{proposition}[\cite{ooKBOMooSkKHvu}]      \label{PROPooNGZJooHjtMyn}
    Soient \( A\in \eM(n,\eK)\), et \( v\in \eK^n\). Si \( B\) est la matrice \( A\) avec la substitution \( L_j\to L_j+v\) et \( C\) est la matrice \( A\) avec la substitution \( L_j\to v\), alors
    \begin{equation}
        \det(B)=\det(A)+\det(C).
    \end{equation}
\end{proposition}

\begin{proof}
    En utilisant l'associativité de la multiplication,
    \begin{subequations}
        \begin{align}
            \det(B)&=\sum_{\sigma\in S_n}\epsilon(\sigma)\prod_{i=1}^nB_{i\sigma(i)}\\
            &=\sum_{\sigma}\epsilon(\sigma)\big( \prod_{i\neq j}B_{i\sigma(i)} \big)B_{j\sigma(j)}\\
            &=\sum_{\sigma}\epsilon(\sigma)\big( \prod_{i\neq j}A_{i\sigma(i)} \big)(A_{j\sigma(j)}+v_{\sigma(j)})\\
            &=\sum_{\sigma}\epsilon(\sigma)\prod_iA_{i\sigma(i)}+\sum_{\sigma}\epsilon(\sigma)\prod_{i\neq j}C_{i\sigma(i)}v_{\sigma(j)}         \label{SUBEQooKATCooVIbEpv}\\
            &=\det(A)+\sum_{\sigma}\epsilon(\sigma)\prod_{i\neq j}C_{i\sigma(i)}C_{j\sigma(j)}  \label{SUBEQooCOTDooPPrEYJ}\\    
            &=\det(A)+\det(C).
        \end{align}
    \end{subequations}
    Justifications :
    \begin{itemize}
        \item \ref{SUBEQooKATCooVIbEpv} parce que pour \( i\neq j\) nous avons \( A_{i\sigma(i)}=C_{i\sigma(i)}\)
        \item \ref{SUBEQooCOTDooPPrEYJ} parce que \( v_{\sigma(j)}=C_{j\sigma(j)}\).
    \end{itemize}
\end{proof}

\begin{proposition}[Combinaison de lignes ou colonnes \( L_k\to L_k+\lambda L_l\)\cite{ooKBOMooSkKHvu}]     \label{PROPooPYNHooLbeVhj}
    Soient une matrice \( A\in \eM(n,\eK)\), deux entiers naturels \( k\neq l\) inférieurs ou égaux à \( n\).
    \begin{enumerate}
        \item       \label{ITEMooJSRDooTggEyO}
            Si \( B\) est la matrice obtenue à partir de \( A\) par la substitution \( L_k\to L_k+\lambda L_l\) ou \( C_k\to C_k+\lambda C_l\), alors
            \begin{equation}
                \det(A)=\det(B).
            \end{equation}
        \item   \label{ITEMooHKZWooVZDgnf}
            Si \( B\) est la matrice \( A\) dans laquelle nous avons fait la substitution \( L_k\to L_k+\lambda L_l\), alors
            \begin{equation}
                B=T_{kl}(\lambda)A
            \end{equation}
            avec \( T_{kl}(\lambda)=\mtu+\lambda E_{kl}\), c'est-à-dire que \( U\) est une matrice de combinaison de lignes.
        \item           \label{ITEMooPGYJooWTTghT}
            La matrice \( T_{kl}(\lambda)\) vérifie \( \det(T_{kl}(\lambda))=1\).
        \item       \label{ITEMooBBEAooZJVGNV}
            Nous avons 
            \begin{equation}
                \det(T_{kl}(\lambda)A)=\det(T_{kl}(\lambda))\det(A).
            \end{equation}
    \end{enumerate}
\end{proposition}

\begin{proof}
    Point par point.
    \begin{subproof}
    \item[\ref{ITEMooJSRDooTggEyO}]
        Soit la matrice \( C\) obtenue à partir de \( A\) par \( L_k\to \lambda L_l\). En considérant le vecteur \( v=\lambda L_l\), nous sommes dans la situation de la proposition \ref{PROPooNGZJooHjtMyn}. Donc
        \begin{equation}
            \det(B)=\det(A)+\det(C).
        \end{equation}
        Mais dans la matrice \( C\), nous avons \( L_k=\lambda L_l\), ce qui implique \( \det(C)=0\) par le corolaire \ref{CORooAZFCooSYINvBl}. Donc \( \det(A)=\det(B)\) comme il se devait.
    \item[\ref{ITEMooHKZWooVZDgnf}]
        Encore un calcul :
        \begin{equation}
            (T_{kl}(\lambda)A)_{ij}=\sum_m\big( \delta_{im}+\lambda(E_{kl})_{im} \big)A_{mj}=A_{ij}+\lambda\sum_m\delta_{ki}\delta_{lm}A_{mj}=A_{ij}+\lambda \delta_{li}A_{kj}.
        \end{equation}
        Cela donne, pour \( i=k\) la ligne
        \begin{equation}
            (T_{kl}(\lambda)A)_{kj}=A_{kj}+\lambda A_{lj},
        \end{equation}
        ce qui correspond bien à \( L_k\to L_k+\lambda L_l\).

    \item[\ref{ITEMooPGYJooWTTghT}]
            Nous calculons le déterminant de \( T_{kl}(\lambda)=\delta+\lambda E_{kl}\) avec \( k\neq l\). Nous avons dans un premier temps :
            \begin{equation}
                \det(T_{kl}(\lambda))=\sum_{\sigma\in S_n}\epsilon(\sigma)\prod_{i=1}^n(\delta_{i\sigma(i)}+\lambda \delta_{ki}\delta_{l\sigma(i)}).
            \end{equation}
 
             Pour les \( \sigma\neq \id\), certains \(\delta_{i\sigma(i)}\) seront nuls; le produit ne sera pas nul si la partie \( \lambda\delta_{ki}\delta_{l\sigma(i)}\) ne s'annule pas. Cela n'arrive qu'à la condition que \( i=k\) et \( \sigma(i)=l\). Donc le seul terme non nul autre que \( \sigma=\id\) peut provenir de \( \sigma=(k,l)\). Pour ce terme, nous isolons les facteurs \( i=l\) et \( i=k\) :
            \begin{equation}
                (\delta_{k\sigma(k)}+\lambda\delta_{kk}\delta_{k\sigma(k)})(\delta_{l\sigma(l)}+\lambda\delta_{kl}\delta_{k\sigma(l)}) =  (\delta_{kl}+\lambda\delta_{kk}\delta_{kl})(\delta_{lk}+\lambda\delta_{kl}\delta_{kl}) = (0 + \lambda 0)(0 + \lambda 0).
            \end{equation}
            Finalement, même avec \( \sigma=(k,l)\), le terme est nul. Il s'ensuit qu'il ne reste \(\sigma = \id \), et
            \begin{equation}
            \det(T_{kl}(\lambda)) = 1 \prod_{i=1}^n(\delta_{ii}+\lambda \delta_{ki}\delta_{li}) = 1
            \end{equation}
            car \(\delta_{ki}\delta_{li} = 0 \) pour tout \( i \).
        \item[\ref{ITEMooBBEAooZJVGNV}]
            En mettant bout à bout les résultats prouvés,
            \begin{equation}
                \det(T_{kl}(\lambda)A)=\det(A)=\det(T_{kl}(\lambda))\det(A).
            \end{equation}
    \end{subproof}
\end{proof}

\begin{proposition}[Multiplication par un scalaire d'une ligne ou colonne \( L_k\to \lambda L_k\)\cite{ooKBOMooSkKHvu}] \label{PROPooXUFKooOaPnna}
    Soient une matrice \( A\in \eM(n,\eK)\), un entier \( k\neq l\) inférieurs ou égal à \( n\). Soit la matrice \( B\) obtenue à partir de \( A\) en multipliant la ligne \( L_k\) par \( \lambda\in \eK\). Alors:
    \begin{enumerate}
        \item       \label{ITEMooBKIGooCDQEDt}
            \( \det(B)=\lambda\det(A)\)
        \item       \label{ITEMooWRRCooFXkRNW}
            En considérant la matrice \( D_k(\lambda)=\mtu+(\lambda-1)E_{kk}\), nous avons
            \begin{equation}
                B=D_k(\lambda)A,
            \end{equation}
            c'est-à-dire que la matrice \( D_k(\lambda) \) est une matrice de multiplication de ligne par un scalaire.
        \item       \label{ITEMooOGGDooPVVRzk}
            Nous avons \( \det(D_k(\lambda))=\lambda\).
        \item       \label{ITEMooIFRVooWQYgkK}
            Et aussi : \( \det(D_k(\lambda)A)=\det(D_k(\lambda))\det(A)\)
    \end{enumerate}
\end{proposition}

\begin{proof}
    Point par point.
    \begin{subproof}
        \item[\ref{ITEMooBKIGooCDQEDt}]
            La matrice \( B\) est donnée par les éléments
            \begin{equation}
                B_{ij}=\begin{cases}
                    A_{ij}    &   \text{si } j\neq k\\
                    \lambda A_{ij}    &    \text{si } j=k
                \end{cases}
            \end{equation}
            c'est-à-dire \( B_{ij}=\big( 1+(\lambda-1)\delta_{jk} \big)A_{ij}\). Nous mettons cela dans la définition du déterminant de \( B\) :
            \begin{equation}        \label{EQooGVMTooPntKew}
                \det(B)=\sum_{\sigma\in S_n}\epsilon(\sigma)\prod_{i=1}^nB_{i\sigma(i)}=\sum_{\sigma}\prod_i\big( 1+(\lambda-1)\delta_{\sigma(i)k}A_{i\sigma(i)} \big).
            \end{equation}
            L'associativité du produit dans \( \eK\) nous permet de séparer le produit de la façon suivante :
            \begin{equation}
                \prod_{i=1}^n\big( 1+(\lambda-1)\delta_{\sigma(i)k} \big)A_{i\sigma(i)}=\prod_i\big( 1+(\lambda-1)\delta_{\sigma(i)k} \big)\prod_iA_{i\sigma(i)}=\lambda\prod_iA_{i\sigma(i)}.
            \end{equation}
            En remettant dans \eqref{EQooGVMTooPntKew}, nous trouvons \( \det(B)=\det(A)\).
        \item[\ref{ITEMooWRRCooFXkRNW}]
            C'est un cas particulier de la proposition \ref{PROPooPYNHooLbeVhj}\ref{ITEMooHKZWooVZDgnf} en prenant \( k=l\) et en adaptant le \( \lambda\).
        \item[\ref{ITEMooOGGDooPVVRzk}]
            Nous calculons le déterminant de la matrice \( D_k(\lambda)=\mtu+(\lambda-1)E_{kk}\). La formule du déterminant donne
            \begin{equation}
                \det(D_k(\lambda))=\sum_{\sigma}\epsilon(\sigma)\prod_{i=1}^n\big( \delta_{i\sigma(i)}+(\lambda-1)\delta_{ki}\delta_{k\sigma(i)} \big).
            \end{equation}
            Si \( i\neq \sigma(i)\), alors non seulement \( \delta_{i\sigma(i)}=0\), mais en plus \( \delta_{ki}\delta_{k\sigma(i)}=0\). Donc seul \( \sigma=\id\) reste dans la somme sur \( \sigma\in S_n\). Il reste donc
            \begin{equation}
                \det(D_k(\lambda))=\prod_{i=1}^n\big( 1+(\lambda-1)\delta_{ki} \big)=\left( \prod_{i\neq k}1 \right)(1+(\lambda-1))=\lambda
            \end{equation}
            où nous avons utilisé encore l'associativité pour isoler le facteur \( i=k\).
        \item[\ref{ITEMooIFRVooWQYgkK}]
            Il faut mettre bout à bout les résultats déjà faits :
            \begin{equation}
                \det(D_k(\lambda)A)=\lambda\det(A)=\det(D_k(\lambda))\det(A).
            \end{equation}
    \end{subproof}
\end{proof}

%--------------------------------------------------------------------------------------------------------------------------- 
\subsection{Réduction de Gauss}
%---------------------------------------------------------------------------------------------------------------------------

Nous avons vu les matrices d'opérations élémentaire sur les lignes et colonnes :
\begin{itemize}
    \item Permutation de lignes \( L_k\leftrightarrow L_l\)  : \( P_{kl}=\mtu+E_{kl}+E_{lk}-E_{kk}-E_{ll}\), proposition \ref{PROPooFQRDooRPfuxk}.
    \item Combinaisons de lignes \( L_k\to L_k+\lambda L_l\) : \( T_{kl} (\lambda)=\mtu+\lambda E_{kl}\), proposition \ref{PROPooPYNHooLbeVhj}.
    \item Multiplication d'une ligne par un scalaire \( L_k\to \lambda L_k\) : \( D_k(\lambda) =\mtu+(\lambda-1)E_{kk}\), proposition \ref{PROPooXUFKooOaPnna}.
\end{itemize}

Ces matrices seront dans la suite notées \( G\). Et elles vérifient la grosse propriété
\begin{equation}        \label{EQooLQTVooBYjVYl}
    \det(GA)=\det(G)\det(A)
\end{equation}
pour toute matrice \( A\).

\begin{proposition}[Réduction de Gauss\cite{MonCerveau}]        \label{PROPooJBTZooNLobpf}
    Soit une matrice \( A\in \eM(n,\eK)\) de déterminant non nul : \( \det(A)\neq 0\). Alors il existe des matrices \( G_1,\ldots, G_N\) toutes de type \( P\), \( T \) ou \( D\) telles que
    \begin{equation}
        G_1\ldots G_NA=\mtu.
    \end{equation}
\end{proposition}

\begin{proof}
    Nous faisons une récurrence sur \( n\). D'abord pour \( n=1\), la matrice \( A\) contient un seul élément \( A_{11}\) qui est non nul par hypothèse. Nous pouvons multiplier sa ligne par \( 1/A_{11}\) pour obtenir le résultat. Plus précisément, nous avons l'égalité
    \begin{equation}
        T_{11}(\frac{1}{ A_{11} })A=\delta
    \end{equation}
    dans \( \eM(1, \eK)\). Notons que \( \eK\) est un corps (donc \( A_{11}\) est inversible) commutatif, ce qui permet d'écrire \( 1/A_{11}\) sans ambiguïtés.

    Supposons le résultat prouvé pour \( n\), et voyons ce qu'il se passe pour \( n+1\). Vu que \( \det(A)\neq 0\), aucune de ses colonnes n'est nulle (corolaire \ref{CORooAZFCooSYINvBl}). Il existe donc un \( k\) tel que \( A_{k1}\neq 0\).

    Par la proposition \ref{PROPooFQRDooRPfuxk}, la matrice
    \begin{equation}
        B^{(1)}=P_{k1}A
    \end{equation}
    est une matrice telle que \( B^{(1)}_{11}=A_{k1}\neq 0\). Ensuite, par la proposition \ref{PROPooXUFKooOaPnna} la matrice
    \begin{equation}
        B^{(2)}=D_1(\frac{1}{ A_{k1} })B^{(1)}
    \end{equation}
    vérifie \( B^{(2)}_{11}=1\). 

    Vu que la multiplication par la matrice \( T_{kl}(\lambda)\) fait par la proposition \ref{PROPooPYNHooLbeVhj} la substitution \( L_k\to L_{k}+\lambda L_l\), la matrice
    \begin{equation}
        B^{(3)}=\prod_{k=2}^{n+1}T_{k1}(-B^{(1)}_{k1})B^{(1)}
    \end{equation}
    a toute sa première colonne nulle à l'exception de \( B^{(3)}_{11}=1\).

    Nous n'avons pas donné de nom ni démontré de théorèmes à propos de la substitution \( C_k\to C_k+\lambda C_l\). En passant éventuellement par les transposées et en utilisant les lemmes \ref{LEMooRSJTooQEoOtN} et \ref{LEMooCEQYooYAbctZ} nous obtenons une matrice \( T'{kl}(\lambda)\) ayant la propriété que la matrice
    \begin{equation}
        B^{(4)}=\prod_{k=2}^{n+1}T'{k1}(-B^{(3)}_{1k})B^{(3)}
    \end{equation}
    vérifie \( B^{(4)}_{1j}=B^{(4)}_{j1}=0\) pour tout \( j\) sauf \( j=1\). En d'autres termes, la matrice \( B^{(4)}\) est de la forme
    \begin{equation}
        B^{(4)}=\begin{pmatrix}
            1    &   \begin{matrix} 
                0    &   \ldots    &   0    
            \end{matrix}\\ 
            \begin{matrix}
                0    \\ 
                \vdots    \\ 
                0    
            \end{matrix}&   \begin{pmatrix}
                    &       &       \\
                    &   A'    &       \\
                    &       &   
            \end{pmatrix}
        \end{pmatrix}
    \end{equation}
    où \( A'\) est une matrice de taille \( n\).
    
    Voyons quelques propriétés de \( A'\). Nous savons que
    \begin{equation}
        B^{(4)}=\prod_i G_iA
    \end{equation}
    où les \( G_i\) sont de type \( P\), \( T\) ou \( D\). Vu que \( \det(PA)=\det(P)\det(A)\) (et idem pour \( T\) et \( D\)), nous avons
    \begin{equation}
        \det(B^{(4)})=\prod_i\det(G_i)\det(A),
    \end{equation}
    et comme aucun des \( \det(G_i)\) n'est nul, nous avons encore \( \det(B^{(4)})\neq 0\), ce qui implique \( \det(A')\neq 0\).

    La récurrence peut avoir lieu. Il existe des matrices \( G'_i\) telles que
    \begin{equation}
        G'_1\ldots G'_MA'=\mtu_n
    \end{equation}
    où les \( G'_i\) sont de taille \( n\). En remarquant que
    \begin{equation}
       P_{n+1; k,l} =\begin{pmatrix}
            1    &   \begin{matrix} 
                0    &   \ldots    &   0    
            \end{matrix}\\ 
            \begin{matrix}
                0    \\ 
                \vdots    \\ 
                0    
            \end{matrix}& P_{n; k-1,l-1}
        \end{pmatrix},
    \end{equation}
    et pareillement pour les matrices \( T\) et \( D\), nous voyons qu'en prenant
    \begin{equation}
       G_i =\begin{pmatrix}
            1    &   \begin{matrix} 
                0    &   \ldots    &   0    
            \end{matrix}\\ 
            \begin{matrix}
                0    \\ 
                \vdots    \\ 
                0    
            \end{matrix}& G'_i
        \end{pmatrix},
    \end{equation}
    nous avons
    \begin{equation}
        \prod_{i=1}^MG_iB^{(3)}=
      \begin{pmatrix}
            1    &   \begin{matrix} 
                0    &   \ldots    &   0    
            \end{matrix}\\ 
            \begin{matrix}
                0    \\ 
                \vdots    \\ 
                0    
            \end{matrix}& \prod_{i=1}^MG'_iA'
        \end{pmatrix}=\mtu_{n+1}.
    \end{equation}
\end{proof}

\begin{proposition} \label{PROPooPMYCooAAtHsB}
    Si \( A\in \eM(n,\eK)\) est telle que \( \det(A)=0\), alors il existe des matrices de manipulation de lignes et de colonnes \( G_1,\ldots, G_N\) telles que \( G_1\ldots G_NA\) ait une colonne de zéros.
\end{proposition}

\begin{proof}
    Si la matrice \( A\) elle-même n'a pas de colonnes de zéros, alors nous pouvons faire un pas de réduction de Gauss et obtenir des matrices \( G_1,\ldots,  G_{N_1}\) telles que
    \begin{equation}
        G_1\ldots G_{N_1}A=
      \begin{pmatrix}
            1    &   \begin{matrix} 
                0    &   \ldots    &   0    
            \end{matrix}\\ 
            \begin{matrix}
                0    \\ 
                \vdots    \\ 
                0    
            \end{matrix}& A^{(1)}
        \end{pmatrix}.
    \end{equation}
    Si \( A^{(1)}\) ne possède pas de colonnes de zéros, nous pouvons continuer. 

    Si nous parvenons à faire \( n\) pas de la sorte, alors nous aurions
    \begin{equation}
        G_1\ldots G_NA=\mtu,
    \end{equation}
    et donc \( \det(G_1\ldots G_N)\det(A)=1\), ce qui est impossible lorsque \( \det(A)=0\). Nous en concluons que le processus doit s'arrêter et qu'une des matrices \( A^{(k)}\) doit avoir une colonne de zéros\footnote{En réalité, le processus tel que nous l'avons décrit ne s'arrête que lorsque la première colonne est remplie de zéros.}.
\end{proof}

%--------------------------------------------------------------------------------------------------------------------------- 
\subsection{Matrices inversibles}
%---------------------------------------------------------------------------------------------------------------------------

\begin{propositionDef}      \label{PROPooMLWRooRWfZXE}
    Soit une matrice \( A\in \eM(n,\eK)\). Si les matrices \( B_1\) et \( B_2\) de \( \eM(n,\eK)\) vérifient
    \begin{equation}
        AB_1=B_1A=\mtu
    \end{equation}
    et
    \begin{equation}
        AB_2=B_2A=\mtu,
    \end{equation}
    alors \( B_1=B_2\). Dans ce cas, nous disons que \( A\) est inversible et nous notons \( A^{-1}\) l'unique matrice telle que \( AA^{-1}=A^{-1}A=\mtu\).
\end{propositionDef}

\begin{proof}
    La preuve est réalisée dans le cas général par le lemme \ref{LEMooECDMooCkWxXf}. Mais si vous en voulez une preuve avec les notations d'ici, en voici une.

    Nous avons $AB_1=AB_2$. En multipliant à gauche par \( B_1\), nous trouvons \( B_1AB_1=B_1AB_2\). En remplaçant \( B_1A\) par \( \mtu\) des deux côtés, il reste \( B_1=B_2\).
\end{proof}

\begin{lemma}[\cite{ooKBOMooSkKHvu}]        \label{LEMooGZCTooQigDvC}
    Si \( A\in \eM(n,\eK)\), alors il existe au plus une matrice \( B\in \eM(n,\eK)\) telle que \( AB=\mtu\).
\end{lemma}

\begin{proof}
    Soient des matrices \( B,C\in \eM(n,\eK)\) telles que \( AB=AC=\delta\). Nous allons montrer que \( B=C\).

    Pour cela nous considérons les applications linéaires \( f_A, f_B, f_C\in \End(\eK^)\) associées par la proposition \ref{PROPooGXDBooHfKRrv}. Vu que \( AB=\delta\), par la proposition \ref{PROPooCSJNooEqcmFm}, nous avons \( f_A\circ f_B = f_{AB}=\id\). La proposition \ref{PROPooADESooATJSrH} nous dit alors que \( f_A\) et \( f_B\) sont bijectives. 

    En particulier, vu que \( \{e_i\}\) est une base, son image par \( f_B\) est une base par la proposition \ref{PROPooZFKZooBGLSex}. La proposition \ref{PROPooHLUYooNsDgbn} dit alors que \( \{f_B(e_i)\}\) est une base. Soit \( k \) fixé. Nous décomposons \( f_B(e_k)-f_C(e_k)\) dans cette base :
    \begin{equation}
        f_B(e_k)-f_C(e_k)=\sum_j\alpha^{(k)}_jf_B(e_j).
    \end{equation}
    Vu que \( f_A\circ(f_B-f_C)=0\), nous avons 
    \begin{equation}
        0=f_A\big( f_B(e_k)-f_C(e_k) \big)=\sum_j(f_A\circ f_B)(e_j)=\sum_j\alpha^{(k)}_je_j.
    \end{equation}
    Donc les \( \alpha^{(k)}_j\) sont tous nuls, pour tout \( j \), et quel que soit le \( k \) de dépaet

    Nous en déduisons que \( f_B(e_k)=f_C(e_k)\), et donc \( f_B=f_C\). Cela implique que \( B=C\) par la proposition \ref{PROPooGXDBooHfKRrv}\ref{ITEMooHSMLooRJZref}.
\end{proof}

\begin{proposition}[\cite{ooKBOMooSkKHvu}]      \label{PROPooECIIooVMCIwz}
    Si \( A,B\in \eM(n,\eK)\) vérifient \( AB=\mtu\), alors \( BA=\mtu\).
\end{proposition}

\begin{proof}
    L'astuce est de poser \( C=BA-\delta+B\) et de montrer que \( C=B\). Pour cela, un rapide calcul commence par montrer que
    \begin{equation}
        AC=ABA-A+AB=AB=\delta.
    \end{equation}
    Donc \( C\) est également un inverse à droite de \( A\). Le lemme \ref{LEMooGZCTooQigDvC} donne alors \( C=B\).
\end{proof}

\begin{corollary}       \label{CORooBQLXooTeVfgb}
    Soit \( A\in \eM(n,\eK)\). Si il existe \( B\in \eM(n,\eK)\) tel que \( AB=\mtu\), alors \( A\) est inversible et son inverse est \( B\).
\end{corollary}

\begin{proof}
    Il s'agit d'une paraphrase de la proposition \ref{PROPooECIIooVMCIwz} et de la définition \ref{PROPooMLWRooRWfZXE}.
\end{proof}

\begin{probleme}
  On peut trouver des inverses sans passer par les applications linéaires: le faire, parce que les appli lin, on les voit plus tard.
\end{probleme}

\begin{lemma}       \label{LEMooZDNVooArIXzC}
    Si une matrice \( A\) n'est pas inversible, alors le produit \( AB\) n'est inversible pour aucune matrice \( B\).
\end{lemma}

\begin{proof}
    Supposons que \( AB\) soit inversible. Alors
    \begin{equation}
        AB(AB)^{-1}=\mtu,
    \end{equation}
    ce qui dirait que \( B(AB)^{-1}\) serait un inverse de \( A\).
\end{proof}


%--------------------------------------------------------------------------------------------------------------------------- 
\subsection{Inversibilité et déterminant}
%---------------------------------------------------------------------------------------------------------------------------

\begin{proposition}     \label{PROPooAVIXooMtVCet}
    Une matrice au déterminant non nul est inversible.
\end{proposition}

\begin{proof}
    Si \( A\) est une matrice telle que \( \det(A)\neq 0\), alors la proposition \ref{PROPooJBTZooNLobpf} nous donne des matrices \( G_1,\ldots, G_N\) telles que
    \begin{equation}
        G_1\ldots G_NA=\delta.
    \end{equation}
    Donc la matrice \( G_1\ldots G_N\) est un inverse de \( A\) par le corolaire \ref{CORooBQLXooTeVfgb}.
\end{proof}

\begin{proposition}     \label{PROPooEOKBooKUROFg}
    Si une matrice \( A\) a une ligne ou une colonne de zéros, alors
    \begin{enumerate}
        \item
            \( \det(A)=0\),
        \item
            \( A\) n'est pas inversible.
    \end{enumerate}
\end{proposition}

\begin{proof}
    Par définition nous avons
    \begin{equation}
        \det(A)=\sum_{\sigma\in S_n}\epsilon(\sigma)\prod_{i=1}^nA_{i\sigma(i)}.
    \end{equation}
    Si la \( k\)\ieme ligne est nulle, alors \( A_{k\sigma(k)}=0\) pour tout \( \sigma\). Donc tous les produits contiennent un facteur nul. Donc \( \det(A)=0\).

    Pour toute matrice \( B\) nous avons
    \begin{equation}
        (AB)_{kk}=\sum_lA_{kl}B_{lk}.
    \end{equation}
    Si la \( k\)\ieme ligne de \( A\) est nulle nous avons \( (AB)_{kk}=0\) et donc pas \( AB=\delta\). Donc \( A\) n'est pas inversible.
\end{proof}

\begin{proposition}     \label{PROPooVUDJooLWjmSI}
    Une matrice dont le déterminant est nul n'est pas inversible.
\end{proposition}

\begin{proof}
    Par la proposition \ref{PROPooPMYCooAAtHsB}, il existe des matrices de manipulation de lignes et de colonnes \( G_1,\ldots, G_N\) telles que la matrice \( G_1\ldots G_NA\) ait une colonne de zéros. De là, la proposition \ref{PROPooEOKBooKUROFg} implique que la matrice
    \begin{equation}        \label{EQooQGXBooXxFOtb}
        G_1\ldots G_NA
    \end{equation}
    n'est pas inversible. Vu les déterminants des matrices \( G_i\),  la proposition \ref{PROPooAVIXooMtVCet} implique que \( G_1\ldots G_N\) est inversible. Si \( A\) était inversible, nous aurions
    \begin{equation}
        G_1\dots G_NAA^{-1}(G_1\ldots G_N)^{-1}=\delta,
    \end{equation}
    c'est-à-dire que \( A^{-1}(G_1\ldots G_N)^{-1}\) serait un inverse de la matrice \eqref{EQooQGXBooXxFOtb}. Cette dernière n'ayant pas d'inverse, nous concluons que \( A\) n'en a pas non plus.
\end{proof}

\begin{theorem}     \label{THOooSNXWooSRjleb}
    Une matrice sur un corps commutatif est inversible si et seulement si son déterminant est non nul.
\end{theorem}

\begin{proof}
    Dans un sens c'est la proposition \ref{PROPooAVIXooMtVCet} et dans l'autre sens c'est la proposition \ref{PROPooVUDJooLWjmSI}.
\end{proof}

%---------------------------------------------------------------------------------------------------------------------------
\subsection{Quelques ensembles de matrices particuliers}
%---------------------------------------------------------------------------------------------------------------------------
Certains ensembles de matrices ont une importance particulière, que nous développerons plus tard.

\begin{definition}[Groupe linéaire de matrices]
On note \( \GL(n,A) \) l'ensemble des matrices carrées d'ordre \( n \) à coefficients dans \( A \), qui sont inversibles. En d'autres termes, \( \GL(n,A) = U (\eM (n,A) ) \).
\end{definition}

\begin{definition}[Groupe orthogonal de matrices]\label{DefMatriceOrthogonale}
    On dit qu'une matrice \( M \) est \defe{orthogonale}{matrice!orthogonale} si son inverse est sa transposée, c'est-à-dire si \( M^{-1} = M^t \). On note \( \gO(n, A) \) l'ensemble des matrices carrées d'ordre \( n \) à coefficients dans \( A \), qui sont orthogonales.
\end{definition}

%--------------------------------------------------------------------------------------------------------------------------- 
\subsection{Déterminant et combinaisons de lignes et colonnes}
%---------------------------------------------------------------------------------------------------------------------------
\label{SUBSECooKMSVooBBHwkH}

\begin{proposition}     \label{PROPooUCZVooPkloQp}
    Soient des matrices \( A,B\in \eM(n,\eK)\) telles que \( \det(A)\neq 0\neq \det(B)\). Alors
    \begin{equation}
        \det(AB)=\det(A)\det(B).
    \end{equation}
\end{proposition}

\begin{proof}
    La proposition \ref{PROPooJBTZooNLobpf} nous donne des matrices de permutations de lignes et de colonnes \( G_1,\ldots, G_N\) et \( G'_1,\ldots, G'_N\) telles que\footnote{Les plus acharnés préciseront que pour avoir le même \( N\) des deux côtés, il a fallu compléter avec des matrices \( \delta\) là où il y en avait le moins.}
    \begin{subequations}        \label{EQooDNZUooHBhcZj}
        \begin{align}
            G_1\ldots G_NA&=\delta\\
            G'_1\ldots G'_NB&=\delta.
        \end{align}
    \end{subequations}
    Nous avons
    \begin{equation}
        (G'_1\ldots G'_N)\underbrace{(G_1\ldots G_N)A}_{=\delta}B=\delta.
    \end{equation}
    En prenant le déterminant des deux côtés et en tenant compte de \eqref{EQooLQTVooBYjVYl},
    \begin{equation}
        1=\det(\delta)=\det\big(  G'_1\ldots G'_NG_1\ldots G_NAB\big)=\det(G_1'\ldots G_N')\det(G_1\ldots G_N)\det(AB).
    \end{equation}
    Mais en même temps, les équations \ref{EQooDNZUooHBhcZj} donnent
    \begin{subequations}
        \begin{align}
            \det(G_1\ldots G_N)=\det(A)^{-1}\\
            \det(G_1'\ldots G'_N)=\det(B)^{-1}.
        \end{align}
    \end{subequations}
    Cela pour dire que
    \begin{equation}
        1=\det(A)^{-1}\det(B)^{-1}\det(AB),
    \end{equation}
    et donc ce qu'il nous fallait.
\end{proof}

\begin{proposition}     \label{PROPooWVJFooTmqoec}
    Soient des matrices \( A,B\in \eM(n,\eK)\) telles que \( \det(A)=0\) et \( \det(B)\neq 0\). Alors
    \begin{equation}
        \det(AB)=\det(BA)=\det(A)\det(B).
    \end{equation}
\end{proposition}

\begin{proof}
    Il existe des matrices de manipulations de lignes et de colonnes \( G_1,\ldots, G_N\) telles que \( G_1\ldots G_NB=\delta\). Donc
    \begin{equation}
        0=\det(A)=\det(G_1\ldots G_NBA)=\det(G_1\ldots G_N)\det(BA).
    \end{equation}
    Donc \( \det(BA)=0\).
\end{proof}

\begin{proposition}     \label{PROPooHQNPooIfPEDH}
    Soient des matrices \( A\) et \( B\) sur un corps commutatif. Alors
    \begin{equation}
        \det(AB)=\det(A)\det(B).
    \end{equation}
\end{proposition}

\begin{proof}
    Les propositions \ref{PROPooUCZVooPkloQp} et \ref{PROPooWVJFooTmqoec} ont déjà fait une grosse partie du travail. Il ne reste que le cas où \( \det(A)=\det(B)=0\).

    Dans ce cas, les matrices \( A\) et \( B\) ne sont pas inversibles (proposition \ref{THOooSNXWooSRjleb}). Le produit \( AB\) n'est alors pas inversible non plus\footnote{Citez le lemme \ref{LEMooZDNVooArIXzC} si vous voulez justifier ça.}. La proposition \ref{THOooSNXWooSRjleb}, utilisée dans le sens inverse, nous dit alors que \( \det(AB)=0\).

    Au final dans le cas \( \det(A)=\det(B)=0\) nous avons \( 0=\det(AB)=\det(A)\det(B)=0\).
\end{proof}

Faisons maintenant le cas général des manipulations de lignes et colonnes.

\begin{proposition}     \label{PROPooSLLGooSZjQrv}
    Soit une matrice carrée \( A\in \eM(n,\eK)\). La matrice \( B\) obtenue par la substitution simultanée
    \begin{equation}
        C_j\to \sum_ka_{kj}C_k
    \end{equation}
    a pour déterminant
    \begin{equation}
        \det(B)=\det(a)\det(A).
    \end{equation}
\end{proposition}

\begin{proof}
    L'élément \( B_{ij}\) de la matrice \( B\) est une combinaison linéaire de tous les éléments de sa ligne :
    \begin{equation}
        B_{ij}=\sum_ka_{kj}A_{ik}=(Aa)_{ij}.
    \end{equation}
    Donc \( B=Aa\). La proposition \ref{PROPooHQNPooIfPEDH} nous dit alors que \( \det(B)=\det(a)\det(A)\).
\end{proof}
%---------------------------------------------------------------------------------------------------------------------------
\subsection{Matrices équivalentes et semblables}
%---------------------------------------------------------------------------------------------------------------------------

\begin{definition}  \label{DefBLELooTvlHoB}
    Deux relations d'équivalence entre les matrices.
    \begin{enumerate}
        \item   \label{ItemPFXCooOUbSCt}
    Deux matrices \( A\) et \( B\) sont \defe{équivalentes}{matrice!équivalence} dans \( \eM(n,\eK)\) s'il existe \( P,Q\in\GL(n,\eK)\) telles que \( A=PBQ^{-1}\).
\item
    Deux matrices sont \defe{semblables}{matrices!similitude} s'il existe une matrice \( P\in \GL(n,\eK)\) telle que \( A=PBP^{-1}\).
    \end{enumerate}
\end{definition}

\begin{lemma}   \label{LemZMxxnfM}
    Une matrice de rang\footnote{Définition~\ref{DEFooVVBYooJbliTi}.} \( r\) dans \( \eM(n,\eK)\) est équivalente à la matrice par blocs
    \begin{equation}
        J_r=\begin{pmatrix}
            \mtu_r    &   0    \\
            0    &   0
        \end{pmatrix}.
    \end{equation}
\end{lemma}
\index{rang!classe d'équivalence}

\begin{proof}
    Nous devons prouver que pour toute matrice \( A\in\eM(n,\eK)\) de rang \( r\), il existe \( P,Q\in\GL(n,\eK)\) telles que \(QAP=J_r\). Soit \( \{ e_i \}\) la base canonique de \( \eK^n\), puis \( \{ f_i \}\) une base telle que \( Af_i=0\) dès que \( i>r\), qui existe par le lemme~\ref{LEMVecsaRgFixe}.

    Nous considérons la matrice inversible \( P\) telle que \( Pe_i=f_i\); ses colonnes sont donc précisément les \( f_i \), si bien que
    \begin{equation}
        APe_i=Af_i=\begin{cases}
            0    &   \text{si } i>r\\
            \neq 0    &    \text{sinon}.
        \end{cases}
    \end{equation}
    La matrice \( AP\) se présente donc sous la forme
    \begin{equation}
        AP=\begin{pmatrix}
            M    &   0    \\
            *    &   0
        \end{pmatrix}
    \end{equation}
    où \( M\) est une matrice \( r\times r\). Nous considérons maintenant une base \( \{ g_i \}_{i=1,\ldots, n}\) dont les \( r\) premiers éléments sont les \( r\) premières colonnes de \( AP\) et une matrice inversible \( Q\) telle que \( Qg_i=e_i\). Alors
    \begin{equation}
        QAPe_i=\begin{cases}
            e_i    &   \text{si } i<r\\
            0    &    \text{sinon}.
        \end{cases}.
    \end{equation}
    Cela signifie que \( QAP\) est la matrice \( J_r\).
\end{proof}

\begin{corollary}[Équivalence et rang]      \label{CorGOUYooErfOIe}
    Deux matrices sont équivalentes\footnote{Définition~\ref{DefBLELooTvlHoB}\ref{ItemPFXCooOUbSCt}.} si et seulement si elles sont de même rang.
\end{corollary}

\begin{proof}
    D'abord il y a des implicites dans l'énoncé. Vu que nous voulons soit par hypothèse soit par conclusion que les matrices \( A\) et \( B\) soient équivalentes, nous supposons qu'elles ont même dimension. Soient donc \( A\) et \( B\) deux matrices carrées d'ordre \( n \).

    Par le lemme~\ref{LemZMxxnfM}, deux matrices de même rang \( r\) sont équivalentes à \( J_r\). Elles sont donc équivalentes entre elles.

    Inversement, supposons que \( A\) et \( B\) soient deux matrices équivalentes : \( A=PBQ^{-1}\) avec \( P\) et \( Q\) inversibles. Alors
    \begin{subequations}
        \begin{align}
            \Image(PBQ^{-1})&=\{ PBQ^{-1}v\tq v\in \eK^n \}\\
            &=PB\underbrace{\{ Q^{-1}v\tq v\in \eK^n \}}_{=\eK^n}\\
            &=P\big( B(\eK^n) \big).
        \end{align}
    \end{subequations}
    L'ensemble \( B(\eK^n)\) est un sous-espace vectoriel de \( \eK^n\). Vu que le rang de \( P\) est maximum, la dimension de \( P\big( B(\eK^n) \big)\) est la même que celle de \( B(\eK^n)\). Par conséquent
    \begin{equation}
        \dim\Big( \Image(PBQ^{-1}) \Big)=\dim\big( B(\eK^n) \big)=\rang(B).
    \end{equation}
    Le membre de gauche de cela n'est autre que \( \rang(A)=\dim\big( \Image(PBQ^{-1}) \big)\).
\end{proof}

%---------------------------------------------------------------------------------------------------------------------------
\subsection{Algorithme des facteurs invariants}
%---------------------------------------------------------------------------------------------------------------------------

\begin{proposition}[Algorithme des facteurs invariants\cite{KXjFWKA}]   \label{PropPDfCqee}
    Soit \( (\eA,\delta)\) un anneau euclidien muni de son stathme  et \( U\in \eM(n \times m,\eA)\). Alors il existe \( d_1,\ldots, d_s\in \eA^*\) et des matrices \( P\in\GL(m,\eA)\), \( Q\in \GL(n,\eA)\) tels que nous ayons
    \begin{equation}
        U=P \begin{pmatrix}
            \begin{matrix}
                d_1    &       &       \\
                    &   \ddots    &       \\
                    &       &   d_s
            \end{matrix}&   0    \\
            0    &   0
        \end{pmatrix}Q
    \end{equation}
    avec \( d_i\divides d_{i+1}\) pour tout \( i\).
\end{proposition}
\index{anneau!euclidien!facteurs invariants}
\index{algorithme!facteurs invariants}

\begin{proof}
    Nous allons donner la preuve plus ou moins sous forme d'algorithme.

    D'abord si \( U=0\) c'est bon, on a la réponse. Sinon, nous prenons l'élément \( (i_0,j_0)\) dont le stathme est le plus petit et nous l'amenons en \( (1,1)\) par les permutations
    \begin{equation}
        \begin{aligned}[]
            C_1&\leftrightarrow C_{j_0}\\
            L_1&\leftrightarrow L_{i_0}
        \end{aligned}
    \end{equation}
    Ensuite nous traitons la première colonne jusqu'à amener des zéros partout en dessous de \( u_{11}\) de la façon suivante : pour chaque ligne successivement nous calculons la division euclidienne
    \begin{equation}
        u_{i1}=qu_{11}+r_i,
    \end{equation}
    et nous faisons
    \begin{equation}
        L_i\to L_i-qL_1,
    \end{equation}
    c'est-à-dire que nous enlevons le maximum possible et il reste seulement \( r_i\) en \( u_{i1}\). Vu que le but est de ne laisser que des zéros dans la première colonne, si le reste n'est pas zéro nous ne sommes pas content\footnote{S'il est zéro, nous passons à la ligne suivante}. Dans ce cas nous permutons \( L_1\leftrightarrow L_i\), ce qui aura pour effet de strictement diminuer le stathme de \( u_{11}\) parce qu'on va mettre en \( u_{11}\) le nombre \( r_i\) dont le stathme est strictement plus petit que celui de \( u_{11}\).

    En faisant ce jeu de division euclidienne puis échange, on diminue toujours le stathme de \( u_{11}\), donc ça finit par s'arrêter, c'est-à-dire qu'à un certain moment la division euclidienne de \( u_{i1}\) par \( u_{11}\) va donner un reste zéro et nous serons content.

    Une fois la première colonne ramenée à la forme
    \begin{equation}
        C_1=\begin{pmatrix}
            u_{11}    \\
            0    \\
            \vdots    \\
            0
        \end{pmatrix},
    \end{equation}
    nous faisons tout le même jeu avec la première ligne en faisant maintenant des sommes divisions et permutations de colonnes. Notons que ce faisant nous ne changeons plus la première colonne.

    En fin de compte nous trouvons une matrice\footnote{Nous nommons toujours par la même lettre \( U\) la matrice originale et la modifiée, comme il est d'usage en informatique.}
    \begin{equation}
        U=\begin{pmatrix}
            u_{11}   &   0    &   \ldots    &   0    \\
             0   &       &       &       \\
             \vdots   &       &   A    &       \\
             0   &       &       &
         \end{pmatrix}
    \end{equation}
    Si l'élément \( u_{11}\) ne divise pas un des éléments de \( A\), disons \( a_{ij}\), alors nous faisons
    \begin{equation}
        C_1\to C_1-C_j.
    \end{equation}
    Cela nous détruit un peu la première colonne, mais ne change pas \( u_{11}\). Nous avons maintnant
    \begin{equation}
        U=\begin{pmatrix}
            u_{11}   &   0    &   \ldots    &   0    \\
             0   &       &       &       \\
             *   &       &       &       \\
             u_{ij}   &       &   A    &       \\
             *   &       &       &       \\
             0   &       &       &
         \end{pmatrix}
    \end{equation}
    Et nous refaisons tout le jeu depuis le début. Cependant lorsque nous allons nous attaquer à la ligne \( i\), \( u_{11}\) ne divisera pas \( u_{ij}\), ce qui donnera lieu à une division euclidienne et un échange \( L_1\leftrightarrow L_i\). L'échange consistant à mettre \( r_i\) à la place de \( u_{11}\) et inversement  diminuera encore strictement le stathme. Encore une fois nous allons travailler jusqu'à avoir la matrice sous la forme
    \begin{equation}    \label{EqADcNVgI}
        U=\begin{pmatrix}
            u_{11}   &   0    &   \ldots    &   0    \\
             0   &       &       &       \\
             \vdots   &       &   A    &       \\
             0   &       &       &
         \end{pmatrix},
    \end{equation}
    sauf que cette fois le stathme de \( u_{11}\) est strictement plus petit que la fois précédente. Si \( u_{11}\) ne divise toujours pas tous les éléments de \( A\), nous recommençons encore et encore. En fin de compte nous finissons par avoir une matrice de la forme \eqref{EqADcNVgI} avec \( u_{11}\) qui divise tous les éléments de \( A\).

    Une fois que cela est fait, il faut continuer en recommençant tout sur la matrice \( A\). Nous avons maintenant
    \begin{equation}
        U=\begin{pmatrix}
            \begin{matrix}
                u_{11}  &       \\
                &   u_{22}
            \end{matrix}&   0    \\
            0    &   B
        \end{pmatrix}.
    \end{equation}
    Sous cette forme nous avons \( u_{11}\divides u_{22}\) et \( u_{11}\) divise tous les éléments de \( B\). En effet \( u_{11}\) divisant tous les éléments de \( A\), il divise toutes les combinaisons de ces éléments. Or tout l'algorithme ne consiste qu'à prendre des combinaisons d'éléments.

    Nous finissons donc bien sûr une matrice comme annoncée. De plus n'ayant effectué que des combinaisons de lignes, nous avons seulement multiplié par des matrices inversibles (lemme~\ref{LemyrAXQs}).
\end{proof}
